\chpt{Theoretical foundations of training future specialists using digital history} \label{ch:theory}

\section{Digital history in pedagogical research} \label{sec:dh-review}

Digital history is a modern interdisciplinary approach to the study of history using digital means \cite{ROMEIN, guiliano, Milligan}. There are various interpretations of this term, but they all relate to society, research, teaching, and communication. Digital history involves creating digital archives, databases, and multimedia resources to make historical data more accessible and interactive \citep{Cohen2005}. Unlike traditional historical methodology, digital history does not offer a revolutionary new approach but provides relevant tools in the era of information technologies, enabling the performance of scholarly work in significantly less time \cite{CohenRosenzweig2006, Graham20151, Graham200810}.

The concept of ``digital history'' emerged in the mid-1990s alongside the development of the World Wide Web \cite{SeefeldtClark2009, NCPH2016}. Roy Rosenzweig and Daniel Cohen at the Center for History and New Media (George Mason University) were among the pioneers who articulated the theoretical foundations of this field \cite{CohenRosenzweig2006}. Since then, digital history has expanded to encompass a broad range of activities including text mining, Geographic Information Systems (GIS), digital storytelling, virtual museum exhibits, and computational approaches to historical analysis \cite{Crymble2021, Graham20151, Bordalejo2025, Lucchesi2019, Crymble202185}.

\subsection{Methodology of the systematic review}

To establish the current state of digital history in pedagogical research, we conducted a systematic review following the PRISMA 2020 guidelines \citep{Page2021}. The review aimed to answer the primary research question \cite{Korniienko2025els}: \textit{How effective is teaching digital history to students compared to non-digital methods of teaching history in terms of improving learning outcomes?}

Subsidiary research questions addressed \cite{Korniienko2025els}: (1) the effect of course status (mandatory vs. elective) on effectiveness; (2) which specific digital tools demonstrate the highest effectiveness; (3) which outcome indicators are most sensitive to digital methods; and (4) whether intervention duration affects effectiveness.

\textbf{Search strategy.} The Scopus database was selected as the primary information source due to its comprehensive coverage of peer-reviewed literature across disciplines \cite{scopus}, indexing over 48,000 active titles from more than 5,000 international publishers. The search string employed Boolean operators and was applied to title, abstract, and keyword fields:

\begin{quote}
\texttt{TITLE-ABS-KEY ("digital history" AND student*)}
\end{quote}

The query was further limited to English-language publications and to document types including journal articles, conference papers, and book chapters. The search was executed on July 26, 2025, with no date restrictions applied during the search phase.

\textbf{Inclusion criteria.} Studies were included if they: (1) involved undergraduate or master's-level students enrolled in history courses at accredited higher education institutions, aged 18+; (2) implemented any digital technology intervention designed to enhance history education (collaborative writing platforms, digital archives, GIS, virtual reality, gamification, digital storytelling, computational text analysis) \cite{Wright-Maley2018, Block2021, Tilak2023, Nygren2014, Corlett-Rivera2017, Cheng2023, Koh2013, Ali2023, Joshi2021, Darmawan2025, Hastomo2024}; (3) employed quantitative, qualitative, or mixed methods with empirical data; (4) were published as peer-reviewed articles, conference proceedings, or book chapters between 2011 and 2025.

\textbf{Exclusion criteria.} Studies were excluded if they: (1) involved pre-university (K-12) populations exclusively; (2) were theoretical or conceptual papers lacking empirical data \cite{Wright-Maley2018603, Block2021371, Nygren201478}; (3) focused on technology training for educators without student outcome assessment; (4) addressed disciplines other than history; (5) involved interventions lasting less than two weeks or fewer than five participants; (6) were published in languages other than English.

\textbf{Selection process.} The selection process followed a two-stage screening protocol \cite{Korniienko2025els}. The initial search yielded 46 records from Scopus. After automated deduplication (n=1), 45 unique citations were screened by title and abstract. To support the screening process, Claude~4 Opus and Claude~4 Sonnet were used as auxiliary tools for preliminary analysis, with all final inclusion/exclusion decisions made by human reviewers \cite{MalekiKarami2024, Taneri2025}. At the title-abstract screening stage, 28 records were excluded: 11 involved pre-university populations, 4 examined unrelated interventions, and 13 lacked empirical data.

Of the 17 studies advancing to full-text assessment, 2 could not be retrieved despite attempts to contact authors. The remaining 15 studies were assessed against all inclusion criteria using both human and LLM assistance. Nine additional studies were excluded: 2 were published outside the eligibility timeframe, 3 involved graduate students exclusively, 2 presented insufficient methodological detail, and 2 described theoretical frameworks without empirical testing. Ultimately, 6 studies satisfied all inclusion criteria.

\textbf{Data collection.} Data from each included study were collected by three reviewers simultaneously: two generative AI-based services (Claude~4 Sonnet and Grok~3) and one human researcher \cite{Rubinstein2025, MalekiKarami2024}. The LLMs worked with identical prompts to ensure standardisation. When discrepancies were found, responses from all three reviewers were compared, with final decisions made by the human researcher. A standardised data extraction form captured study characteristics, participant demographics, intervention details, outcome measures, and quality indicators \cite{KitchenhamCharters2007, HigginsThomas2019}.

\subsection{Results of the systematic review}

The identification of only 6 studies meeting the inclusion criteria underscores the nascent state of digital history as a pedagogical research area \cite{Korniienko2025els}. This finding is particularly striking given the widespread use of digital tools in historical research \cite{Graham20151, Crymble2021} and the growing emphasis on digital competencies in education policy \cite{Voogt20161, Vuorikari2016}. All six studies were published in English in peer-reviewed venues between 2011 and 2017 \cite{Bell2016, Coleborne2011, Davis2017, Lee2014159, McLean20144, Soh2013}.

\subsubsection{Characteristics of included studies}

Table~\ref{tab:dh-characteristics} presents the key characteristics of the six included studies, and Table~\ref{tab:dh-interventions} summarises the intervention details.

\begin{table}[htbp]
\centering
\caption{Characteristics of included studies in the digital history systematic review}
\label{tab:dh-characteristics}
\small
\begin{tabular}{p{2.5cm}p{2.5cm}lp{2.2cm}lp{1.5cm}}
\hline
\textbf{Authors} & \textbf{Design} & \textbf{Country} & \textbf{Participants} & \textbf{Level} & \textbf{Duration} \\
\hline
Bell (2016) & Case study & Australia & $\sim$100 & UG & 1 semester \\
Coleborne (2011) & Case study & New Zealand & 43 & UG & 3 weeks \\
Davis (2017) & Case study & USA & Not specified & UG & 1 semester \\
Lee (2014) & Design-based & USA & 200 (total) & UG & 5 years \\
McLean (2014) & Mixed methods & Canada & 27 (8 resp.) & MA & 1 semester \\
Soh (2013) & Quasi-experimental & USA & 150 & UG & 1 semester \\
\hline
\end{tabular}
\end{table}

\begin{table}[htbp]
\centering
\caption{Intervention details for included studies}
\label{tab:dh-interventions}
\small
\begin{tabular}{p{2.5cm}p{3.5cm}p{2.5cm}p{1.5cm}p{2.5cm}}
\hline
\textbf{Authors} & \textbf{Intervention} & \textbf{Technologies} & \textbf{Control} & \textbf{Outcomes} \\
\hline
Bell (2016) & 3-minute digital video stories & Various digital tools & None & Engagement, understanding \\
Coleborne (2011) & Digital storytelling & Windows Movie Maker & None & Storytelling skills, historical thinking \\
Davis (2017) & Digital exhibition & Omeka (web) & None & Motivation, digital skills \\
Lee (2014) & Personal digital histories & Various (iterative) & None & Engagement, technology skills \\
McLean (2014) & Historical projects & Wiki, blogs & None & Technology skills, collaboration \\
Soh (2013) & Collaborative writing & ClassroomWiki & Traditional & Writing skills, engagement, learning \\
\hline
\end{tabular}
\end{table}

The six studies represent diverse geographical contexts: three from the United States \cite{Davis2017, Lee2014159, Soh2013}, and one each from Australia \cite{Bell2016}, New Zealand \cite{Coleborne2011}, and Canada \cite{McLean20144}. Sample sizes ranged from 27 to approximately 200 participants, with the sole quasi-experimental study \cite{Soh2013} having the second-largest sample (150 students). Only one study included a control group, reflecting the common challenge of implementing randomised designs in educational settings \cite{Schulz2010CONSORT}.

\subsubsection{Study findings}

The included studies documented the following interventions and findings:

\begin{enumerate}
\item \textit{Digital storytelling with video} \cite{Bell2016}: Approximately 100 undergraduate students created 3-minute digital video stories about historical topics. The case study documented enhanced engagement with historical content and deeper understanding of historical concepts through the process of narrative construction. The digital format facilitated students' spontaneous pursuit of additional primary sources beyond course requirements.

\item \textit{Digital storytelling with multimedia} \cite{Coleborne2011}: Forty-three undergraduates used Windows Movie Maker to create digital stories about historical events. The study found improvements in both digital storytelling skills and historical thinking, demonstrating that multimedia production can serve as a vehicle for developing analytical capabilities.

\item \textit{Digital exhibition creation} \cite{Davis2017}: Students used the Omeka web platform to create public-facing digital exhibitions incorporating metadata standards and Dublin Core elements. The study reported that 45\% of students expressed enthusiasm about the digital format, and students demonstrated improved digital skills alongside historical learning.

\item \textit{Personal digital histories} \cite{Lee2014159}: This design-based research study spanned 5 years with 8 iterative cycles, involving approximately 200 students. Students created personal digital histories using various technologies. The iterative refinement process revealed that student engagement was highest when digital tools were directly aligned with historical inquiry tasks \cite{Reeves2003}.

\item \textit{Wiki and blog-based projects} \cite{McLean20144}: Twenty-seven master's students (8 of whom provided survey responses) used wiki and blog platforms for collaborative historical research projects. The mixed-methods study found improvements in technology skills and collaborative competencies, though the small response rate limited generalisability.

\item \textit{Collaborative wiki writing} \cite{Soh2013}: The only quasi-experimental study, involving 150 undergraduate students using ClassroomWiki compared to traditional approaches. The study documented a 340\% increase in peer-to-peer interactions compared to traditional discussion formats, with students spending an average of 12.4 hours on digital assignments (approximately 40\% more than on traditional assignments). Writing skills, engagement, and learning outcomes all showed positive effects.
\end{enumerate}

\subsection{Cross-cutting analysis of outcomes}

Analysis of primary and secondary outcomes across the six studies revealed consistent patterns.

\textbf{Primary outcomes.} Historical knowledge acquisition was assessed in four studies, with three (67\%) reporting positive effects. \citet{Soh2013} provided the strongest evidence, with students in the wiki-based condition demonstrating significantly better writing quality and content accuracy than the control group. \citet{Lee2014159} documented improved engagement with historical evidence across eight iterative cycles. \citet{Bell2016} and \citet{Coleborne2011} reported enhanced understanding of historical concepts through the digital storytelling process, though without quantitative comparison to control conditions.

\textbf{Secondary outcomes.} Student engagement was the most consistently positive outcome, with all six studies (100\%) reporting positive effects on at least one dimension of engagement \cite{Bell2016, Coleborne2011, Davis2017, Lee2014159, McLean20144, Soh2013}. Behavioural engagement was measured through platform analytics (time-on-task, number of contributions, peer interactions) and showed the most dramatic improvements. Cognitive engagement manifested through enhanced depth of historical analysis: \citet{Bell2016} reported that students spontaneously pursued additional primary sources beyond course requirements. Source analysis skills improved in three studies \cite{Davis2017, Lee2014159, McLean20144}, particularly those employing digital archive platforms. The development of digital skills emerged as an unexpected but valuable outcome, with students acquiring competencies in metadata creation, digital curation, and platform-specific technical skills \cite{Davis2017, McLean20144}.

\textbf{Mechanisms of effectiveness.} Three mechanisms appeared to drive positive outcomes across studies \cite{Korniienko2025els}: (1) \textit{multimodal representation}, enabling students to engage with history through diverse sensory channels \cite{Bell2016, Coleborne2011}; (2) the \textit{``authentic audience'' effect}, where knowing their digital products would be publicly visible motivated higher-quality work \cite{Davis2017, Lee2014159}; and (3) \textit{technological affordances} that enabled pedagogical approaches impractical in analog environments (version control making thinking visible, collaborative platforms distributing cognitive load) \cite{Soh2013, McLean20144}.

\textbf{Barriers and challenges.} Despite positive outcomes, implementation challenges were documented \cite{Korniienko2025els}. Technical difficulties consumed substantial instructional time (one study reported 15\% of sessions devoted to troubleshooting) \cite{McLean20144}. Instructor preparedness was a critical barrier, with faculty struggling to integrate technological and pedagogical knowledge effectively \cite{Davis2017, Lee2014159}. Assessment challenges were also noted: traditional evaluation methods proved inadequate for capturing learning in digital environments \cite{Coleborne2011}.

\subsection{Analysis and implications}

The small number of identified studies reveals a significant gap between the practice of digital history in research and its systematic integration into pedagogy. Several factors may explain this gap:

First, digital history has developed primarily as a research methodology rather than a teaching approach \cite{Graham20151, Crymble2021}. The pioneers of the field (Rosenzweig, Cohen) focused on demonstrating the utility of digital tools for historical investigation \cite{CohenRosenzweig2006}, with less attention to how these tools can be effectively taught to students. The six included studies all represent individual instructors' initiatives rather than systematic programme-level implementations.

Second, the interdisciplinary nature of digital history creates challenges for pedagogical research \cite{Bordalejo2025, Lucchesi2019, Luther20106571}. Studies must address competencies in both history and technology, requiring expertise that spans traditional disciplinary boundaries. The TPACK framework, referenced in four of the six studies \cite{Bell2016, Davis2017, Lee2014159, McLean20144}, provides a theoretical lens for understanding this intersection but does not resolve the practical challenge of designing interventions that develop both types of knowledge.

Third, the rapid evolution of digital tools means that pedagogical approaches may become obsolete before they can be rigorously evaluated. This creates a tension between innovation and evidence-based practice. The included studies span 2011--2017, using platforms (ClassroomWiki, Windows Movie Maker, Omeka) that have evolved substantially since publication \cite{Zielezinski20141137}.

Fourth, the absence of standardised outcome measures for digital history pedagogy makes it difficult to compare results across studies or to accumulate evidence systematically \cite{Hattie2009, Tamim2011}. Unlike STEM education, where standardised tests provide common outcome metrics, history education relies on diverse assessment approaches (essays, projects, portfolios) that resist direct comparison \cite{MonteSano2012, Reisman2012, Manfra2008223, Sebring2012229}.

The scarcity of rigorous pedagogical research on digital history motivates the broader examination of gamification in history education presented in the following sections, as gamification represents the most extensively researched digital approach to history teaching \cite{Connolly2012, DichevDicheva2017, KoivistoHamari2019}. While the digital history review identified only 6 studies, the gamification review \cite{KorniienkoSemerikov2024} identified 74 -- a 12-fold difference that underscores the disproportionate research attention that game-based approaches have received.

\subsection{Conclusions to Section 1.1}

The PRISMA-guided systematic review \cite{Page2021, Korniienko2025els} demonstrates that digital history in pedagogical research remains an underdeveloped field, with only 6 studies meeting rigorous inclusion criteria \cite{Bell2016, Coleborne2011, Davis2017, Lee2014159, McLean20144, Soh2013}. This finding establishes the need for: (1) more empirical research on digital history pedagogy; (2) examination of specific digital approaches (such as gamification) that have been more extensively studied \cite{Connolly2012, Qian2016}; and (3) development of comprehensive teaching methodologies that integrate multiple digital tools \cite{Voogt20161}.

\section{Gamification in history education: definitions, effectiveness, and design principles} \label{sec:gamification-review}

\subsection{Background and rationale}

Gamification -- the application of game design elements to non-game contexts \cite{Deterding2011} -- has attracted significant research attention in education over the past decade \cite{Hamari2014, Sailer2017, DichevDicheva2017, KoivistoHamari2019}. In history education specifically, gamification offers opportunities to enhance engagement, motivation, and learning through interactive elements such as points, badges, leaderboards, quests, and narrative structures \cite{Qian2016, Connolly2012, Prensky2001}.

Given the paucity of research on digital history pedagogy identified in Section~\ref{sec:dh-review}, this section presents a comprehensive systematic review of gamification in history education \cite{KorniienkoSemerikov2024} as the most extensively researched component of digital history methodology.

\subsection{Methodology}

The systematic review followed the PRISMA 2020 guidelines \citep{Page2021}. Four research questions guided the review:
\begin{enumerate}
\item How is gamification defined and conceptualised in history education?
\item What is the current state of research on gamification effectiveness in history education?
\item What are the key design principles, strategies, and technologies used?
\item What are the main challenges, limitations, and potential negative effects?
\end{enumerate}

\textbf{Search strategy.} Three databases were searched: Scopus \cite{scopus}, Web of Science \cite{wos}, and Dimensions \cite{dimensions}. The search strategy combined gamification-related terms with history education terms in paper titles:

\begin{quote}
\texttt{TITLE (history AND (game* OR gamification) AND (education OR teaching OR learning OR student* OR instruction))}
\end{quote}

Limits were applied for time frame (2000--2024), language (English or Ukrainian), and publication type (peer-reviewed articles, conference papers, book chapters).

\textbf{Inclusion criteria:} (1) students in primary, secondary, or higher education learning history; (2) gamification techniques, elements, or strategies in history education; (3) empirical studies using quantitative, qualitative, or mixed methods; (4) formal or informal educational settings; (5) published between 2000 and 2024; (6) indexed in Scopus, Web of Science, or Dimensions.

\textbf{Screening process.} The initial search yielded 4,011 records across the three databases \cite{KorniienkoSemerikov2024}. After duplicate removal, title-abstract screening, and full-text assessment, 80 studies advanced to data extraction. Of these, 6 were excluded during detailed analysis (3 were not about history education, 1 was a systematic review, 1 was about the history of mathematics, and 1 contained no gamification intervention), yielding 74 studies published between 2004 and 2024 across 27 countries.

Data extraction was automated using GPT-4o with a structured 21-item prompt, followed by human verification \cite{Rubinstein2025, MalekiKarami2024}. Risk-of-bias assessment employed an adapted ROBIS tool with three domains and six signalling questions \cite{Whiting2016, HigginsThomas2019}.

\subsection{Definitional analysis}

\subsubsection{Prevalence of definitional ambiguity}

A critical finding of this review concerns the use of definitions. Of the 74 included studies:

\begin{itemize}
\item Only 11 studies (14.9\%) provided a clear, explicit definition of gamification;
\item 63 studies (85.1\%) used the term without providing an operational definition;
\item Among those providing definitions, the most commonly cited was that of Deterding et al. (2011): ``the use of game design elements in non-game contexts'' \citep{Deterding2011};
\item Several studies conflated gamification with game-based learning, serious games, or educational games without distinguishing between these related but distinct concepts \cite{36Scholz2021, 42Huffer201581, 43ali2018indonesian}.
\end{itemize}

This definitional ambiguity has significant implications for research synthesis and practice \cite{DichevDicheva2017, KoivistoHamari2019, encyclopedia3040089}. Without clear definitions, it is difficult to compare findings across studies, replicate research, or develop evidence-based guidelines for practitioners \cite{Clark2016, Buckley2018}.

The definitional approaches differed systematically by disciplinary orientation. Studies rooted in educational technology and human--computer interaction tended to define gamification narrowly as the addition of game elements to non-game contexts, emphasising mechanics such as points, leaderboards, and badges \cite{Deterding2011, SEABORN201514}. In contrast, studies from history education and the humanities conceptualised the approach more broadly as game-based or experiential learning, foregrounding narrative immersion, historical empathy, and perspective-taking \cite{12Gilbert2019108, 55Squire2004, 56Akkerman2009449}. A third group of studies from computer science and design framed the concept through a software engineering lens, emphasising system design and usability \cite{Souza2018201, Striuk_2022}. This disciplinary fragmentation helps explain both the terminological inconsistency and the diversity of theoretical frameworks employed \cite{Buckley2018, Calderon2016}.

Table~\ref{tab:explicit-definitions} presents the 11 studies that provided explicit definitions. These definitions span a range of conceptualisations, from broad frameworks for game-based learning to specific definitions of gamification, serious games, pervasive games, alternate reality games, and simulation games. The diversity in terminology and scope underscores the field's lack of definitional consensus.

\begin{table}[htbp]
\centering
\caption{Studies providing explicit definitions of gamification or related terms (n=11 of 74)}
\label{tab:explicit-definitions}
\small
\begin{tabular}{p{3cm}p{2.5cm}p{7cm}}
\hline
\textbf{Study} & \textbf{Term defined} & \textbf{Definition (excerpt)} \\
\hline
\citet{30Zin2008} & Game-based learning & ``A paradigm which utilizes the game as a medium for conveying the learning contents; leveraging the power of computer games to captivate and engage end users'' \\
\citet{23Zin2009322} & DGBL & ``A paradigm which utilizes the game as a medium for conveying the learning contents\ldots for a specific purpose'' \\
\citet{76Corbeil2011462} & Simulation game & ``Places [the participant] in a context, with certain givens, where he must attempt to achieve a goal'' \\
\citet{65Razak201139} & Educational game & ``A melding of educational content, learning principles, and computer games\ldots'' \\
\citet{42Huffer201581} & Serious game & ``Games that have an explicit and carefully thought-out educational purpose and are not intended to be played primarily for amusement'' \\
\citet{58Shih2015173} & Pervasive game & ``A mobile game that seamlessly integrates the physical environment with virtual scenarios'' \\
\citet{71lynch2015pedagogical} & Alternate reality game & ``Cross media narrative-based games that use the Internet as a central communications platform'' \\
\citet{43ali2018indonesian} & Gamification & ``A method of applying game mechanics to engage and motivate people more to complete their task voluntarily'' \\
\citet{36Scholz2021} & Gamification vs. GBL & ``Gamification referring to the addition of game-based elements\ldots game-based learning [as] more than isolated game elements'' \\
\citet{34Serrano2023} & Gamification & ``A technique, a method and a strategy at the same time\ldots starts from the knowledge of the elements that make games attractive'' \\
\citet{41Petersen2023301} & Serious game & ``Applications that are not designed exclusively for fun or that are intended to be more than entertainment'' \\
\hline
\end{tabular}
\end{table}

\subsubsection{Conceptual approaches}

We identified four distinct conceptual approaches to gamification in the literature:

\begin{enumerate}
\item \textbf{Structural gamification}: Adding game elements (points, badges, leaderboards) to existing educational content without modifying the content itself \cite{Deterding2011, SEABORN201514}.
\item \textbf{Content gamification}: Transforming educational content to make it more game-like (narratives, characters, challenges, choices) \cite{43ali2018indonesian, 34Serrano2023}.
\item \textbf{Full gamification}: Comprehensive integration of game mechanics, dynamics, and aesthetics into the learning experience \cite{36Scholz2021, Sailer2017}.
\item \textbf{Game-based learning}: Using complete games (digital or analogue) as the primary pedagogical tool, which some studies classified under gamification \cite{30Zin2008, 23Zin2009322, Qian2016, Prensky2001}.
\end{enumerate}

\subsection{Characteristics of the evidence base}

Before examining effectiveness findings, it is important to understand the characteristics of the 74 included studies, as these context factors influence the interpretation of results.

\subsubsection{Geographic distribution}

Research was conducted across 27 countries, but production was highly concentrated. Table~\ref{tab:country-distribution} presents the complete distribution.

\begin{table}[htbp]
\centering
\caption{Distribution of studies by country of application (n=74; some studies report multiple countries)}
\label{tab:country-distribution}
\small
\begin{tabular}{lclc}
\hline
\textbf{Country} & \textbf{n} & \textbf{Country} & \textbf{n} \\
\hline
Taiwan & 13 & United States & 11 \\
Malaysia & 9 & Indonesia & 8 \\
Canada & 6 & Greece & 5 \\
Spain & 4 & United Kingdom & 3 \\
Sweden & 3 & Italy & 3 \\
Poland & 3 & Peru & 2 \\
Netherlands & 2 & Pakistan & 1 \\
Japan & 1 & Tunisia & 1 \\
Saudi Arabia & 1 & Chile & 1 \\
Portugal & 1 & Turkey & 1 \\
Denmark & 1 & Australia & 1 \\
China (Hong Kong) & 1 & Ireland & 1 \\
Austria & 1 & Germany & 1 \\
China & 1 & & \\
\hline
\end{tabular}
\end{table}

Taiwan produced the most studies (n=13, 17.6\%) \cite{16Chen201447, 19Chou2019, 28Chao201734, 44Shiue2016, 52Lin2011265, 58Shih2015173, 64Hsu201324, 67Lin201296, 70Hou201519, 9Lin2018, 18Shih201054, 45Yu2016727, 8Yu2014}, followed by the United States (n=11, 14.9\%) \cite{55Squire2004, 54Gibson2007, 15Slater2022, 12Gilbert2019108, 29Stiso20201165, 57Huntley2011567, 59Kee2014270, 60Hoy2018, 74Fendt2019, 78Rouse2022357, 80gottlieb2018your}, Malaysia (n=9, 12.2\%) \cite{30Zin2008, 23Zin2009322, 43ali2018indonesian, 48Moketar2021, 50Lye2013, 65Razak201139, 69Wong2017150, 2Mahamarowi2023133, 21Ghulamani2017}, and Indonesia (n=8, 10.8\%) \cite{20Erandaru2014, 66Rahmawati2023163, 77Ramansyah2021, 22Nugroho2020149, 72hutauruk2020utilization, 75trajano2014using, 79Fernes2024379, 13Humpire2022}. Together, these four countries accounted for 55.4\% of the total. East and Southeast Asian countries (Taiwan, Malaysia, Indonesia, China/Hong Kong, Japan) contributed 43.2\% of studies, reflecting strong regional interest in gamified history education. European countries contributed 33.8\% of studies, with Greece (n=5) \cite{1Barbatsis201159, 25Repantis2011502, 39Tegos2014838, 51Ardito2007, 5kazanidis2018alpha}, Spain (n=4) \cite{17Manzanares2019, 34Serrano2023, 35Fernandez2023, 53Perez2014}, and the United Kingdom, Sweden, Italy, and Poland (n=3 each) representing the most active contexts. The Global South (Latin America, Africa, Middle East, South Asia) was almost entirely absent, with only 6 studies from these regions combined. Notably, no studies from Ukraine or other post-Soviet states were identified, despite the region's rich historical traditions and growing digital education infrastructure \cite{UNESCO2023}.

\subsubsection{Temporal distribution and research evolution}

Studies spanned 2004--2024 with a pronounced increase from 2011 onward. Table~\ref{tab:yearly-distribution} presents the yearly distribution.

\begin{table}[htbp]
\centering
\caption{Yearly distribution of included studies (n=74)}
\label{tab:yearly-distribution}
\small
\begin{tabular}{cccccccc}
\hline
\textbf{Year} & \textbf{n} & \textbf{Year} & \textbf{n} & \textbf{Year} & \textbf{n} & \textbf{Year} & \textbf{n} \\
\hline
2004 & 1 & 2010 & 2 & 2016 & 1 & 2022 & 4 \\
2007 & 3 & 2011 & 7 & 2017 & 4 & 2023 & 7 \\
2008 & 1 & 2012 & 1 & 2018 & 8 & 2024 & 3 \\
2009 & 3 & 2013 & 2 & 2019 & 6 & & \\
 & & 2014 & 7 & 2020 & 5 & & \\
 & & 2015 & 5 & 2021 & 4 & & \\
\hline
\end{tabular}
\end{table}

The most productive years were 2011, 2014, 2018, and 2023 (each contributing 7--8 studies). The earlier period (2004--2017) and later period (2018--2024) each contributed 37 studies. Analysis of keyword evolution through bibliometric network analysis revealed three distinct temporal stages:

\begin{enumerate}
\item \textbf{Early exploration stage (2004--2012)}: Characterised by initial investigations into using computer games for history education \cite{55Squire2004, 54Gibson2007, 7Natvig2004F-1, 49Huizenga2007127, 30Zin2008, 51Ardito2007}. Keywords such as \textit{computer game}, \textit{educational game}, and \textit{learning} dominated this period. Research focused on foundational questions about whether games could effectively support history learning \cite{Prensky2001, Petr2007629}.

\item \textbf{Narrative and engagement stage (2013--2017)}: A shift toward examining how game narratives and storytelling can enhance historical understanding \cite{12Gilbert2019108, 59Kee2014270, 42Huffer201581, 4Cruz2017, 68cole2015end}. Keywords like \textit{history learning}, \textit{video game}, \textit{motivation}, and \textit{serious game} reflected growing interest in leveraging narrative elements for historical empathy and engagement.

\item \textbf{Immersive technology stage (2018--2024)}: The adoption of emerging technologies \cite{10Huaman2019220, 6Czerwinski2018, 31DeAmicis2009451, 32Cheng20241718, 38Mercik2024, 40Kuran2018, Alexov2023}, with keywords such as \textit{augmented reality}, \textit{virtual reality}, \textit{gamification}, and \textit{mobile game} indicating a broadening of technological approaches and a more systematic integration of game design principles into history education.
\end{enumerate}

\subsubsection{Publication types and database coverage}

The majority of included studies were conference papers (n=40, 54.1\%), followed by journal articles (n=24, 32.4\%), book chapters (n=9, 12.2\%), and one preprint (1.4\%). Scopus was the most comprehensive source, covering 62 of 74 studies (83.8\%). Dimensions contributed 11 unique studies (14.8\%) not indexed in Scopus or Web of Science, demonstrating the value of searching multiple databases. The substantial proportion of conference papers suggests that much of the research in this area is preliminary, with findings presented at conferences before (or instead of) undergoing the more rigorous peer-review process of journal publication \cite{KorniienkoSemerikov2024}.

\subsubsection{Education levels}

Secondary education was the most studied context (n=35, 47.3\%), followed by higher education (n=21, 28.4\%) and primary education (n=13, 17.6\%). Five studies (6.8\%) addressed multiple levels or did not specify. The predominance of secondary education studies may reflect the alignment between game-based approaches and the developmental stage of adolescent learners, who are typically more receptive to gamified instruction \cite{49Huizenga2007127, 11Ternblad2023115, 15Slater2022, Adams2022119}. The relative scarcity of higher education studies is notable given that university students represent the target population for NMT preparatory courses \cite{Ahlfeld2021493}.

\subsubsection{Game types and technological formats}

The interventions employed diverse formats, as shown in Table~\ref{tab:game-types}.

\begin{table}[htbp]
\centering
\caption{Distribution of game types across included studies (n=74)}
\label{tab:game-types}
\small
\begin{tabular}{lcc}
\hline
\textbf{Game type} & \textbf{n} & \textbf{\%} \\
\hline
VR/AR games & 19 & 25.7 \\
Board and card games & 11 & 14.9 \\
Custom digital games & 10 & 13.5 \\
Web-based games & 10 & 13.5 \\
Commercial off-the-shelf & 8 & 10.8 \\
Other/mixed & 7 & 9.5 \\
Mobile games & 5 & 6.8 \\
RPG & 3 & 4.1 \\
Simulation & 1 & 1.4 \\
\hline
\end{tabular}
\end{table}

VR/AR games were the most common (25.7\%), reflecting recent technological adoption \cite{1Barbatsis201159, 10Huaman2019220, 31DeAmicis2009451, 6Czerwinski2018, 73Wolfartsberger2011160}. Board and card games (14.9\%) \cite{13Humpire2022, 40Kuran2018, 46Stirling2022, 76Corbeil2011462, 38Mercik2024} and custom digital games (13.5\%) \cite{16Chen201447, 19Chou2019, 50Lye2013, 64Hsu201324, 9Lin2018} were also frequently used. Eight studies (10.8\%) employed commercial off-the-shelf games such as \textit{Civilization IV} \cite{55Squire2004} or \textit{Assassin's Creed} \cite{12Gilbert2019108}. Digital formats accounted for approximately 84\% of the corpus.

The temporal evolution of game types reveals important trends. VR/AR technologies have been present since 2007 \cite{51Ardito2007} but show markedly increased adoption from 2018 onward \cite{6Czerwinski2018, 10Huaman2019220, 32Cheng20241718}, coinciding with the decreased cost of VR headsets and the improved accessibility of AR through smartphones. Custom digital games have been consistently popular throughout the period. Board and card games emerged as a notable category from 2014 \cite{59Kee2014270, 46Stirling2022, 40Kuran2018}, reflecting growing interest in non-digital gamification approaches that avoid technological barriers. Commercial off-the-shelf titles were sporadically used across the period, with notable examples including the Civilization series \cite{55Squire2004, 57Huntley2011567, 60Hoy2018, 80gottlieb2018your}, Assassin's Creed \cite{12Gilbert2019108, 78Rouse2022357}, and various historical strategy games \cite{frith2021, Barboutidis2023613}.

\subsubsection{Research designs}

The design landscape was dominated by non-experimental approaches. Table~\ref{tab:research-designs} presents the distribution of designs.

\begin{table}[htbp]
\centering
\caption{Distribution of research designs across included studies (n=74)}
\label{tab:research-designs}
\small
\begin{tabular}{lcc}
\hline
\textbf{Design} & \textbf{n} & \textbf{\%} \\
\hline
Qualitative & 14 & 18.9 \\
Quasi-experimental & 13 & 17.6 \\
Descriptive & 10 & 13.5 \\
Experimental (true) & 9 & 12.2 \\
Mixed methods & 9 & 12.2 \\
Design-based research & 7 & 9.5 \\
Pre-post (single group) & 4 & 5.4 \\
Case study & 3 & 4.1 \\
Survey & 2 & 2.7 \\
Other & 3 & 4.1 \\
\hline
\end{tabular}
\end{table}

The experimental family accounted for 36.5\% (n=27), but only 1 study employed a fully randomised controlled design \cite{Schulz2010CONSORT}. Quasi-experimental designs (n=13) \cite{8Yu2014, 52Lin2011265, 24Ghannem2019, 49Huizenga2007127} and pre-experimental single-group pretest-posttest designs (n=9) \cite{1Barbatsis201159, 19Chou2019, 9Lin2018, 63Liu2021} made up the rest. Qualitative and case study designs constituted 21.6\% (n=16) \cite{12Gilbert2019108, 55Squire2004, 78Rouse2022357}, design-based research 13.5\% (n=10), and mixed methods 12.2\% (n=9). Only 23 of 74 studies (31.1\%) included a control or comparison group, and a substantial number (n=43, 58.1\%) employed no quantitative comparative methodology at all. This methodological landscape stands in stark contrast to fields such as medical education or psychology, where randomised controlled trials are the expected standard for effectiveness claims \cite{Schulz2010CONSORT, GuyattGRADE2008, SterneEtAl2016}.

\subsubsection{Sample sizes and statistical power}

Among the 62 studies reporting participant numbers, the median was 33.5 (mean = 127.0, range: 4--2,417). The distribution was heavily right-skewed, as shown in Table~\ref{tab:sample-sizes}.

\begin{table}[htbp]
\centering
\caption{Distribution of sample sizes across included studies (n=74)}
\label{tab:sample-sizes}
\small
\begin{tabular}{lcc}
\hline
\textbf{Sample size range} & \textbf{n} & \textbf{\%} \\
\hline
$<$30 & 25 & 33.8 \\
30--99 & 23 & 31.1 \\
100--499 & 9 & 12.2 \\
$\geq$500 & 5 & 6.8 \\
Not reported & 12 & 16.2 \\
\hline
\end{tabular}
\end{table}

The prevalence of small samples (40.3\% below 30, 66.1\% below 50 among those reporting) has direct implications for statistical power \cite{Cohen1960}. A two-group comparison with medium effect size (d = 0.50) and $\alpha$ = .05 requires approximately 64 participants per group (128 total) to achieve 80\% power, a threshold exceeded by only a minority of studies. Studies with particularly small samples included \citet{1Barbatsis201159} (n=17), \citet{22Nugroho2020149} (n=5), \citet{51Ardito2007} (n=4), and \citet{73Wolfartsberger2011160} (n=4). These small samples substantially limit generalisability and increase the likelihood of both Type~I errors (due to inflated variance) and Type~II errors (due to insufficient power). Furthermore, publication bias assessment through funnel plot analysis \cite{Sterne2011} was not feasible given the heterogeneity of outcome measures.

\subsubsection{Theoretical frameworks}

The studies cited diverse theoretical frameworks, with game-based learning frameworks being the most frequently referenced (27 studies) \cite{30Zin2008, 23Zin2009322, Prensky2001, Qian2016}, followed by flow theory (11 studies) \cite{16Chen201447, 19Chou2019, 70Hou201519}, experiential learning and constructivism (9 studies each) \cite{Constructivism, Bonwell1991}, collaborative learning and situated learning (8 studies each). Motivation theories were cited in 7 studies \cite{Sailer2017, Hamari2014}, as were historical thinking frameworks and design frameworks. The Technology Acceptance Model (TAM) appeared in 3 studies, and self-determination theory in 2.

Notably, discipline-specific frameworks such as historical thinking and historical empathy were cited in only 7 studies \cite{12Gilbert2019108, 56Akkerman2009449, 55Squire2004, 78Rouse2022357}, suggesting that gamification research in history education draws more heavily from general educational and psychological theories than from the epistemological traditions of the history discipline itself. This finding has implications for the theoretical grounding of gamification studies: without engagement with what it means to think historically \cite{MonteSano2012, Reisman2012}, gamification interventions risk developing surface-level knowledge (dates, names, events) rather than the deeper analytical skills that characterise expert historical reasoning.

\subsubsection{Comparison with gamification in other domains}

The characteristics of the history education evidence base can be contextualised by comparison with gamification research in other domains. A parallel systematic review of gamification in software engineering education \citep{Korniienko2025aredu} identified 29 studies from an initial pool of 486 records. Several patterns are consistent across domains:

\begin{itemize}
\item \textbf{Positive but uncertain effects}: Both reviews found predominantly positive effects on engagement and motivation, with more mixed evidence for learning outcomes \cite{Korniienko2025aredu, DichevDicheva2017, KoivistoHamari2019}. In software engineering, the average effect on learning outcomes was small (Cohen's d = 0.23), with low certainty of evidence according to GRADE assessment \cite{GRADE, GuyattGRADE2008, Guyatt2008}.
\item \textbf{Methodological limitations}: Both domains are characterised by small sample sizes, lack of control groups, and reliance on unvalidated instruments \cite{Clark2016, Connolly2012}. In software engineering, 18 of 29 studies (62.1\%) had serious overall risk of bias \cite{Korniienko2025aredu}.
\item \textbf{Game element preferences}: Software engineering gamification more commonly employed structural elements (points in 17/29 studies, leaderboards in 11/29, badges in 9/29) \cite{Souza2018201, Striuk_2022, Polat_2023, SEABORN201514, Akpolat2014149, Berkling2013, DetelingRiedl2015, Mora2016, Iosup2014, Hasan2018, Matsubara2017, Bartel2016, Fuchs2016}, while history education relied more heavily on narrative and immersive elements (narrative/storytelling in 44/74 studies, immersion/VR/AR in 55/74).
\item \textbf{Definitional consistency}: Neither domain showed strong definitional practices, though the history education review found even lower rates of explicit definition (14.9\% vs. approximately 30\% in software engineering) \cite{Korniienko2025aredu}.
\end{itemize}

These cross-domain comparisons suggest that the methodological challenges identified in history education gamification research are not unique but rather reflect broader weaknesses in the gamification evidence base across educational domains \cite{DichevDicheva2017, KoivistoHamari2019, Hamari2014, Sailer2017, Morales-Trujillo2021}.

\subsection{Effectiveness findings}

\subsubsection{Outcome types measured}

The review synthesised effectiveness evidence across multiple outcome categories \cite{KorniienkoSemerikov2024}. The most frequently measured outcome was student engagement (52 studies, 70.3\%), followed by knowledge acquisition (45 studies, 60.8\%), attitudes and satisfaction (39 studies, 52.7\%), and motivation (26 studies, 35.1\%). Higher-order thinking skills, including historical thinking and empathy, were assessed in only 14 studies (18.9\%) \cite{12Gilbert2019108, 56Akkerman2009449, 55Squire2004, 78Rouse2022357, 59Kee2014270}, indicating a gap in the evaluation of deeper cognitive and affective outcomes specific to the history discipline. Flow and immersion were measured in 15 studies \cite{16Chen201447, 19Chou2019, 70Hou201519, 28Chao201734}, collaboration in 18, and usability in 13.

\subsubsection{Overall direction of effects}

Of the 74 studies, 12 (16.2\%) reported exclusively positive results across all measured outcomes, 37 (50.0\%) reported mixed results (both positive and null or negative findings across different outcome measures), 18 (24.3\%) reported null or negative results, and 7 (9.5\%) did not measure learning or engagement outcomes quantitatively \cite{KorniienkoSemerikov2024}. The high proportion of mixed results underscores the complexity of gamification effects and the importance of examining specific outcome domains rather than drawing blanket conclusions about ``effectiveness'' \cite{DichevDicheva2017, Clark2016}.

\subsubsection{Student engagement and motivation}

The majority of studies (n=52, 70.3\%) reported positive effects of gamification on student engagement or motivation. Common findings included increased participation in voluntary activities, longer time-on-task, and higher self-reported motivation. Behavioural engagement metrics (platform analytics, participation rates) showed the most consistent positive effects. Specific quantitative findings included:

\begin{itemize}
\item A significant correlation between perceived reality and engagement (r = .70, p $<$ .001) in a VR-based application \cite{1Barbatsis201159};
\item 73\% of students positively motivated in a collaborative gaming study, with 93\% having prior experience with gaming \cite{3Watson2011466};
\item 100\% of surveyed students expressing willingness to play the game again in one mobile-based study \cite{58Shih2015173};
\item Mean engagement scores above 5 on a 7-point scale across multiple VR studies \cite{10Huaman2019220, 31DeAmicis2009451};
\item Over 90\% of students showing positive motivation toward a game-based system, with satisfaction scores averaging 4.51 out of 5 \cite{64Hsu201324};
\item A significant gender difference in engagement, with boys demonstrating statistically higher engagement compared to girls (F(1,64) = 4.408, p = .040) \cite{49Huizenga2007127}.
\end{itemize}

However, the sustainability of these effects beyond initial novelty periods was rarely examined \cite{DichevDicheva2017, Hamari2014}. Some studies documented that not all gamification strategies were equally motivating: competitive timed games generated stress in some students \cite{11Ternblad2023115}, approximately 10\% of participants in one study reported being negatively motivated \cite{47Clara2010156}, and female students were less familiar with or interested in strategy-based games \cite{49Huizenga2007127}. These findings highlight the importance of considering individual differences and offering multiple pathways within gamified environments \cite{Sailer2017}.

\subsubsection{Learning outcomes}

Results regarding academic achievement were more mixed \cite{KorniienkoSemerikov2024}:
\begin{itemize}
\item 38 studies (51.4\%) reported positive effects on learning outcomes;
\item 28 studies (37.8\%) reported no significant difference compared to control conditions;
\item 8 studies (10.8\%) reported negative or mixed effects \cite{47Clara2010156, Oceja2022419}.
\end{itemize}

Table~\ref{tab:significant-learning} presents a selection of studies reporting statistically significant quantitative learning outcomes. Among the studies reporting effect sizes, values ranged from small (d = 0.28) to medium (d = 0.61), with stronger effects observed for factual recall than for interpretive understanding.

\begin{table}[htbp]
\centering
\caption{Selected studies reporting statistically significant learning outcomes}
\label{tab:significant-learning}
\small
\begin{tabular}{p{3cm}clp{3cm}p{3.5cm}}
\hline
\textbf{Study} & \textbf{n} & \textbf{Design} & \textbf{Outcome} & \textbf{Key statistic} \\
\hline
\citet{1Barbatsis201159} & 17 & Pre-post & Knowledge & Z = $-$2.226, p = .026 \\
\citet{8Yu2014} & 103 & Quasi-exp. & Achievement & t = 2.01, p = .05 \\
\citet{9Lin2018} & 27 & Pre-post & Effectiveness & t = 8.196, p $<$ .001 \\
\citet{16Chen201447} & 38 & Pre-post & Knowledge & t = $-$5.33, p $<$ .001 \\
\citet{19Chou2019} & 66 & Pre-post & Knowledge & t = 9.58, p $<$ .001 \\
\citet{24Ghannem2019} & 36 & Quasi-exp. & Knowledge & W = 0, p $\leq$ .01 \\
\citet{52Lin2011265} & 30 & Quasi-exp. & Achievement & Low achievers p = .010 \\
\citet{63Liu2021} & 12 & Pre-post & Knowledge & Z = $-$2.77, p = .006 \\
\hline
\end{tabular}
\end{table}

This pattern suggests that gamification may be more effective for developing surface-level knowledge than for fostering deep historical thinking \cite{Clark2016, Connolly2012}. Qualitative studies provided additional evidence of learning enhancement, including increased correct responses in historical knowledge assessments \cite{24Ghannem2019, 19Chou2019}, higher retention scores in game-based learning compared to traditional methods (up to 9/10 vs. 3/10 correct answers) \cite{43ali2018indonesian}, increased student engagement and critical thinking \cite{12Gilbert2019108, 78Rouse2022357}, improved team collaboration \cite{3Watson2011466, 46Stirling2022}, and the development of systemic understanding of history through recursive cycles of failure and strategy revision \cite{55Squire2004, 57Huntley2011567}.

\subsubsection{Flow and immersion}

Eight studies specifically measured flow states using validated instruments \cite{16Chen201447, 19Chou2019, 28Chao201734, 44Shiue2016, 45Yu2016727, 70Hou201519, 64Hsu201324, 67Lin201296}, providing evidence that history games can achieve meaningful immersion. All reported flow scores above scale midpoints. Concentration and immersion consistently scored highest among flow dimensions, while loss of self-consciousness and time transformation scored lowest \cite{16Chen201447, 70Hou201519}. One study found that flow state was significantly positively correlated with learning achievement (p $<$ .05) \cite{16Chen201447}, providing evidence that achieving flow states in history games may be a mechanism through which engagement translates into learning. Mean flow scores ranged from 3.63 to 4.23 on 5-point scales, with autotelic experience (enjoyment of the activity for its own sake) scoring particularly high (M = 4.54 in one study) \cite{70Hou201519}.

\subsubsection{Historical thinking skills}

A smaller subset of studies (n=18) specifically examined the development of historical thinking skills such as source analysis, contextualisation, and perspective-taking \cite{MonteSano2012, Reisman2012}. Of these, 12 reported positive outcomes, while 6 found no significant effects. Studies employing multi-perspectival digital archives appeared most effective, as students encountered conflicting accounts requiring critical evaluation \cite{59Kee2014270, 60Hoy2018}. Narrative video games were conceptualised as tools for fostering empathy and emotional connections to historical events and figures through role-play elements \cite{12Gilbert2019108, 78Rouse2022357, 56Akkerman2009449}. Board games were viewed as mechanisms for fostering historical perspective-taking through role-playing and strategic elements \cite{46Stirling2022, 76Corbeil2011462, 40Kuran2018}. The limited number of studies assessing higher-order thinking skills (only 18.9\% of the corpus) represents a significant gap, as these skills are central to the goals of history education \cite{MonteSano2012, Reisman2012}.

\subsubsection{Affective outcomes and attitudes toward history}

Studies examining affective outcomes such as attitude toward history, enjoyment, and self-efficacy generally reported positive results (n=31 of 37 reporting on these outcomes, 83.8\%). The strongest positive effects were observed in studies using narrative-rich game environments, where historical content was embedded in engaging storylines \cite{12Gilbert2019108, 74Fendt2019, 29Stiso20201165}. One large-scale study (N = 2,417) identified five distinct engagement profiles through latent class analysis, finding that modifying game dialogue scripts shifted some students toward more positive evaluations of history learning \cite{15Slater2022}. Other studies documented transformations in students' relationships with history, with some students moving from viewing it as ``boring'' to developing genuine interest through game-based exploration \cite{55Squire2004, 80gottlieb2018your}. Structural equation modelling in one study confirmed that perceived usefulness had the strongest influence on game preference, and more sophisticated epistemological beliefs positively affected attitudes toward game-based history learning \cite{45Yu2016727}.

\subsection{Implementation challenges}

The review identified several categories of challenges that recurred across studies:

\textbf{Technical constraints.} Infrastructure requirements, software compatibility, and device availability posed persistent challenges \cite{KorniienkoSemerikov2024}. VR/AR studies (25.7\% of the corpus) were particularly affected by hardware costs and maintenance demands \cite{1Barbatsis201159, 6Czerwinski2018, 73Wolfartsberger2011160}. Web-based and mobile formats showed lower technical barriers but introduced their own issues (browser compatibility, screen size limitations, connectivity requirements) \cite{58Shih2015173, 48Moketar2021}.

\textbf{Resource limitations.} Developing custom gamification environments required substantial investment in design, programming, and content creation \cite{50Lye2013, 20Erandaru2014, 9Lin2018, Gasca-Hurtado2016}. Studies using commercial off-the-shelf titles (10.8\%) avoided development costs but sacrificed pedagogical control \cite{55Squire2004, 12Gilbert2019108, Castronovo2015}. The resource intensity of gamification development may explain the prevalence of small-scale studies: few research teams have the resources to develop, deploy, and evaluate gamification at scale.

\textbf{Curriculum integration difficulties.} Fitting gamification into existing curricula and assessment structures was cited as a significant challenge \cite{61Karsenti202027, 57Huntley2011567}. Time pressure was a common concern: gamified activities typically required more class time than traditional instruction, creating tensions with content coverage requirements \cite{49Huizenga2007127, 46Stirling2022}.

\textbf{Teacher preparation.} Instructor preparedness and confidence with gamification tools was identified as a critical factor \cite{Qian2016, UNESCO2023}. Studies where instructors received substantial training or co-designed the gamification elements reported better outcomes than those where technology was introduced without adequate pedagogical support \cite{61Karsenti202027, 36Scholz2021, Buisman2014}.

\textbf{Student resistance and equity.} Some studies documented initial resistance from students who perceived gamification as trivialising serious historical content \cite{47Clara2010156, Oceja2022419, 33Oceja2022419}, or who were uncomfortable with competitive elements \cite{11Ternblad2023115, Hof2017}. Equity concerns were raised regarding differential access to technology and varying levels of digital literacy among students \cite{UNESCO2023, Voogt20161, Figg20181215}.

\subsection{Game design elements and technological platforms}

\subsubsection{Game design element taxonomy}

Analysis of the 74 studies identified a comprehensive taxonomy of game design elements utilised in history education gamification. Table~\ref{tab:game-elements} presents the frequency of major categories.

\begin{table}[htbp]
\centering
\caption{Frequency of game design element categories across included studies (n=74)}
\label{tab:game-elements}
\small
\begin{tabular}{lcc}
\hline
\textbf{Game design element category} & \textbf{n} & \textbf{\%} \\
\hline
Immersion/VR/AR & 55 & 74.3 \\
Narrative/Storytelling & 44 & 59.5 \\
Role-playing & 44 & 59.5 \\
Collaboration/Competition & 36 & 48.6 \\
Levels/Progression & 35 & 47.3 \\
Points/Rewards/Scoring & 31 & 41.9 \\
Exploration & 29 & 39.2 \\
Quiz/Assessment & 19 & 25.7 \\
Puzzles/Challenges & 18 & 24.3 \\
Rules/Goals & 18 & 24.3 \\
Feedback & 17 & 23.0 \\
Time pressure & 12 & 16.2 \\
\hline
\end{tabular}
\end{table}

Immersive technologies (VR/AR and 3D environments) were employed in 55 studies (74.3\%), making them the most prevalent design element category \cite{1Barbatsis201159, 10Huaman2019220, 6Czerwinski2018, 31DeAmicis2009451, 73Wolfartsberger2011160, 32Cheng20241718}. Role-playing and narrative/storytelling elements were each present in 44 studies (59.5\%) \cite{18Shih201054, 55Squire2004, 12Gilbert2019108, 76Corbeil2011462, 74Fendt2019}, reflecting the natural alignment between historical content and narrative-driven game mechanics. Collaboration and competition mechanics appeared in 36 studies (48.6\%) \cite{3Watson2011466, 49Huizenga2007127, 46Stirling2022, 39Tegos2014838}, while points, rewards, and scoring systems -- often considered the hallmark of gamification \cite{Deterding2011, SEABORN201514} -- were present in only 31 studies (41.9\%). This finding suggests that history education gamification relies more heavily on narrative and experiential elements than on extrinsic reward systems, a pattern that distinguishes it from gamification in other domains (where points and leaderboards typically dominate) \cite{Korniienko2025aredu, Souza2018201, SEABORN201514}.

These elements can be further organised into functional clusters:

\textbf{Exploration and navigation mechanics} encompassed 3D exploration and interactive object manipulation \cite{1Barbatsis201159, 51Ardito2007}, virtual environment exploration \cite{10Huaman2019220, 73Wolfartsberger2011160}, location-based exploration and GPS-triggered events \cite{58Shih2015173, 25Repantis2011502, Colteli2015}, open world exploration with narrative choice and branching \cite{55Squire2004, 74Fendt2019}, first-person and third-person perspective navigation, and map-based navigation. These mechanics leveraged spatial cognition and discovery-based learning to engage students with historical environments.

\textbf{Assessment and feedback systems} included quizzes with increasing difficulty and immediate feedback \cite{19Chou2019, 24Ghannem2019, 34Serrano2023}, multiple-choice questions integrated into gameplay, progressive task difficulty and challenge systems \cite{64Hsu201324, 9Lin2018}, performance tracking and progress monitoring, and real-time feedback through game mechanics \cite{11Ternblad2023115, 15Slater2022}. The integration of assessment within gameplay (rather than as separate evaluative events) was a distinguishing feature of effective implementations.

\textbf{Narrative and storytelling elements} comprised immersive storytelling and digital narratives \cite{12Gilbert2019108, 74Fendt2019, 29Stiso20201165}, character-driven plots and NPC (non-player character) interactions \cite{18Shih201054, 50Lye2013}, historical scenario simulations and role-playing \cite{76Corbeil2011462, 55Squire2004}, in-game cutscenes and cinematics, and storyline progression and chapter-based narratives \cite{80gottlieb2018your, 78Rouse2022357}. Historical narratives were adapted into game stories, real historical events served as plot bases, and historical figures were cast as interactive characters \cite{59Kee2014270, 68cole2015end}.

\textbf{Competition and collaboration mechanisms} included team-based collaboration and cooperative gameplay \cite{3Watson2011466, 49Huizenga2007127, Gomes2014}, competitive elements and leaderboards \cite{34Serrano2023, 43ali2018indonesian, DeSousaPinto2017}, group competition and tournaments \cite{46Stirling2022, 39Tegos2014838, Diniz2017}, and multiplayer functionality. The balance between competition and collaboration was identified as a critical design decision, as purely competitive structures risked discouraging lower-performing students \cite{11Ternblad2023115, Sailer2017}.

\textbf{Reward and progression systems} encompassed point-based scoring and rewards \cite{43ali2018indonesian, 64Hsu201324}, progression-based unlocking of content and level systems \cite{19Chou2019, 9Lin2018}, achievement records and badges \cite{34Serrano2023, Deterding2011}, and lives/health systems with resource management \cite{16Chen201447, 18Shih201054}.

\textbf{Interactive historical content} included historical artifact examination and collection \cite{51Ardito2007, 59Kee2014270}, interactive timelines and concept maps \cite{62Milosz201896, 27Sumi2022}, character cards and historical figure representations \cite{13Humpire2022, 40Kuran2018}, period-specific item cards \cite{46Stirling2022, 38Mercik2024}, and historical riddles and puzzles \cite{35Fernandez2023, 53Perez2014}.

\textbf{Strategic and decision-making elements} comprised strategy-based mechanics and tactical gameplay \cite{55Squire2004, 57Huntley2011567, 80gottlieb2018your}, resource management and economic systems \cite{60Hoy2018}, diplomatic relations and alliance systems, military conflict and combat systems \cite{76Corbeil2011462}, and historical decision-making scenarios \cite{56Akkerman2009449, 47Clara2010156}.

\subsubsection{Technological platforms}

The studies employed a diverse range of technological platforms \cite{KorniienkoSemerikov2024}. Unity 3D/Unity Engine was the most frequently utilised game engine, appearing in 11 studies \cite{1Barbatsis201159, 10Huaman2019220, 6Czerwinski2018, 32Cheng20241718}, followed by Unreal Engine 4 (3 studies) \cite{17Manzanares2019}. Web-based technologies (HTML5, CSS3, JavaScript) were used in browser-based games \cite{48Moketar2021, 5kazanidis2018alpha, 39Tegos2014838}, with PHP and MySQL for backend data management. VR/AR implementations ranged from low-cost solutions (Google Cardboard with smartphones) to high-fidelity systems (HTC Vive) \cite{73Wolfartsberger2011160, 31DeAmicis2009451}. Commercial games repurposed for education included the Civilization series (6 studies) \cite{55Squire2004, 57Huntley2011567, 60Hoy2018, 80gottlieb2018your}, Assassin's Creed franchise (2 studies) \cite{12Gilbert2019108, 78Rouse2022357}, and various complex historical strategy games (Crusader Kings II, Europa Universalis IV, Hearts of Iron IV) \cite{37Lewis202054}. Non-digital platforms -- physical board games, card-based games, and 3D printed models \cite{46Stirling2022, 40Kuran2018, 13Humpire2022, 76Corbeil2011462, 38Mercik2024} -- represented an important subset, demonstrating that gamification does not require digital technology.

Communication and collaboration platforms used alongside gamification included Zoom for online learning, social media (Facebook) for community building, and learning management systems (Moodle) \cite{Gupta2022341, Bozic2024538}. This multi-platform approach foreshadows the multi-channel design adopted in the present study (Chapter~\ref{ch:methodology}).

\subsubsection{Instructional design strategies}

Collaborative learning was the most frequently employed instructional strategy (46 studies, 62.2\%) \cite{3Watson2011466, 49Huizenga2007127, 46Stirling2022}, followed by narrative-driven instruction (40 studies, 54.1\%) \cite{12Gilbert2019108, 74Fendt2019, 29Stiso20201165} and assessment-integrated approaches (31 studies, 41.9\%) \cite{19Chou2019, 24Ghannem2019, 34Serrano2023}. Scaffolded/guided approaches and problem/inquiry-based learning were each present in 23 studies (31.1\%), as was experiential/situated learning \cite{Bonwell1991, Constructivism}. Constructivist approaches, while theoretically prominent in the literature, were explicitly employed in only 15 studies (20.3\%), suggesting a gap between theoretical framing and actual instructional design \cite{Constructivism}.

Content integration approaches included immersive historical reconstruction (virtual reconstruction of archaeological sites) \cite{1Barbatsis201159, 51Ardito2007, 31DeAmicis2009451}, narrative-driven learning (content delivered through storytelling) \cite{12Gilbert2019108, 74Fendt2019}, problem-based learning (historical challenges requiring application of knowledge) \cite{59Kee2014270, 42Huffer201581}, simulation-based learning (games modelling historical systems) \cite{55Squire2004, 76Corbeil2011462}, role-playing and perspective-taking (students assuming roles of historical figures) \cite{18Shih201054, 56Akkerman2009449}, artifact-based learning (interaction with digitised historical artifacts) \cite{51Ardito2007, 62Milosz201896}, location-based learning (GPS-triggered digital content at historical sites) \cite{58Shih2015173, 25Repantis2011502}, and assessment-integrated learning (knowledge questions embedded as game challenges) \cite{24Ghannem2019, 19Chou2019, 43ali2018indonesian}.

\subsection{Design principles}

From the synthesis of effective interventions, six evidence-based design principles for gamification in history education were identified. These principles emerged from the analysis of studies reporting positive outcomes and represent the common features of successful implementations.

\begin{enumerate}
\item \textbf{Historical authenticity}: The most effective implementations maintained historical accuracy while incorporating game elements \cite{12Gilbert2019108, 1Barbatsis201159, 51Ardito2007}. Superficial gamification that trivialised historical content was associated with negative outcomes in 10.8\% of studies \cite{47Clara2010156, Oceja2022419}. Authenticity encompasses both factual accuracy and appropriate representation of historical complexity, ambiguity, and multiple perspectives \cite{MonteSano2012, Reisman2012}. Studies using historically accurate game environments (reconstructions of archaeological sites, period-appropriate visual design) reported stronger learning outcomes than those using generic game formats with historical content superimposed \cite{1Barbatsis201159, 31DeAmicis2009451, 62Milosz201896}.

\item \textbf{Narrative integration}: Using historical narratives as the framework for gamified activities enhanced both engagement and learning \cite{12Gilbert2019108, 74Fendt2019, 29Stiso20201165}. Quest-based structures aligned naturally with historical investigation, casting students as detectives uncovering historical evidence \cite{59Kee2014270, 42Huffer201581}. Narrative-rich environments showed the strongest effects on affective outcomes (attitude toward history, enjoyment) \cite{15Slater2022, 55Squire2004}. Commercial titles like Assassin's Creed demonstrated the potential of narrative integration but raised concerns about historical liberties taken for entertainment purposes \cite{12Gilbert2019108, 78Rouse2022357}.

\item \textbf{Progressive difficulty}: Scaffolded challenges that increased in complexity over time supported sustained engagement and skill development \cite{19Chou2019, 9Lin2018, 64Hsu201324}. Effective scaffolding moved from recognition tasks (matching, multiple choice) to analysis tasks (source evaluation, perspective-taking) to synthesis tasks (constructing historical narratives, arguing causation). This progression mirrors established taxonomies of historical thinking skills \cite{MonteSano2012, Reisman2012, Calandra2005323} and ensures that gamification supports, rather than undermines, cognitive development.

\item \textbf{Immediate feedback}: Providing instant feedback on quiz responses and game actions facilitated self-regulated learning \cite{11Ternblad2023115, 15Slater2022, 24Ghannem2019}. Studies documented that immediate feedback enabled iterative learning cycles where students corrected misconceptions in real time \cite{19Chou2019, 34Serrano2023}. The feedback mechanism was particularly effective for factual recall tasks but less clearly effective for complex interpretive questions where multiple valid answers exist \cite{Clark2016}.

\item \textbf{Social elements}: Collaborative and competitive elements (team challenges, leaderboards, peer review) enhanced motivation, but their effects were moderated by implementation design \cite{3Watson2011466, 49Huizenga2007127, Sailer2017}. Competition needed to be carefully balanced to avoid discouraging lower-performing students \cite{11Ternblad2023115, 47Clara2010156}. Collaborative structures (team-based quests, cooperative problem-solving) showed more consistently positive effects than purely competitive elements \cite{46Stirling2022, 39Tegos2014838}.

\item \textbf{Multiple modalities}: Combining text, images, maps, interactive elements, audio, and video addressed different learning preferences and supported deeper understanding \cite{10Huaman2019220, 62Milosz201896, 27Sumi2022}. The use of primary source images, historical maps, and audiovisual materials within gamified environments was associated with enhanced engagement across diverse learner profiles \cite{59Kee2014270, 74Fendt2019}. VR/AR implementations provided the most immersive multimodal experiences but at higher cost and technical complexity \cite{6Czerwinski2018, 32Cheng20241718, 73Wolfartsberger2011160, Hopf2024375}.
\end{enumerate}

\subsection{Conclusions to Section 1.2}

The systematic review of 74 studies on gamification in history education \cite{KorniienkoSemerikov2024}, drawn from an initial pool of 4,011 records across three databases, reveals a rapidly growing but methodologically immature field. Key findings include:

\begin{itemize}
\item \textbf{Definitional ambiguity}: 85.1\% of studies lack explicit definitions of gamification, undermining construct validity and cross-study comparison;
\item \textbf{Positive engagement effects}: 70.3\% of studies report positive effects on engagement, but with limited evidence on long-term sustainability;
\item \textbf{Mixed learning outcomes}: 51.4\% report positive effects on learning, but effect sizes are small to medium (d = 0.28--0.61) and concentrated on factual recall rather than deep understanding;
\item \textbf{Design diversity}: VR/AR (25.7\%), board games (14.9\%), custom digital (13.5\%), and web-based (13.5\%) formats dominate, with digital approaches accounting for 84\% of the corpus;
\item \textbf{Geographic concentration}: Four countries (Taiwan, USA, Malaysia, Indonesia) account for 55.4\% of all studies;
\item \textbf{Methodological weaknesses}: Only 31.1\% include control groups, median sample size is 33.5, and only 1 study is an RCT.
\end{itemize}

The six design principles identified (historical authenticity, narrative integration, progressive difficulty, immediate feedback, social elements, multiple modalities) provide an evidence-based foundation for implementation, which informs the experimental design presented in Chapter~\ref{ch:methodology}.


\section{Methodological quality of gamification research in history education} \label{sec:mrqs}

\subsection{Rationale for quality assessment}

The positive findings reported in Section~\ref{sec:gamification-review} must be interpreted in light of the methodological quality of the underlying research \cite{DichevDicheva2017, KoivistoHamari2019, Clark2016}. Previous meta-analyses in educational technology have shown that studies with higher methodological quality tend to report smaller effect sizes, suggesting that methodological weaknesses may inflate perceived effectiveness \citep{Hattie2009} \cite{Tamim2011, Connolly2012}.

This section presents a quality assessment of the 74 studies identified in the systematic review \cite{KorniienkoSemerikov2024}, using the Mixed Methods Research Quality Synthesis (MRQS) tool \cite{Korniienko2025mrqs} with an innovative LLM validation approach.

\subsection{The MRQS framework}

To address the absence of a structured quality assessment instrument adapted for gamification research in history education, we developed the Methodological Reporting Quality Score (MRQS) \cite{Korniienko2025mrqs} -- a ten-item binary instrument (scored 0 or 1 per item, range 0--10) designed to assess the completeness and rigour of reporting practices. The MRQS captures fundamental aspects of methodological reporting endorsed in educational research guidelines \cite{Schulz2010CONSORT, GuyattGRADE2008, HigginsThomas2019, Campbell2020} and codable from published study reports.

The ten MRQS items are:

\begin{enumerate}
\item \textbf{M1: Sample size explicitly reported} -- transparent reporting of participant numbers as a minimum requirement for evaluating statistical conclusions \cite{Schulz2010CONSORT};
\item \textbf{M2: Sample size $\geq$ 30} -- a commonly regarded minimum for basic parametric inference \cite{Cohen1960};
\item \textbf{M3: Control or comparison group present} -- necessary for causal inference about intervention effects \cite{Schulz2010CONSORT, HigginsThomas2019};
\item \textbf{M4: Validated measurement instruments used} -- prerequisites for trustworthy outcome assessment;
\item \textbf{M5: Effect sizes or standardised measures reported} -- enabling cross-study comparison, recommended by APA and CONSORT guidelines \cite{Schulz2010CONSORT, Borenstein2009, Higgins2011};
\item \textbf{M6: Pre-post or longitudinal design} -- permitting assessment of within-participant change;
\item \textbf{M7: Inferential statistics reported} -- providing probabilistic evidence regarding the reliability of observed differences;
\item \textbf{M8: Explicit definition of gamification provided} -- necessary for construct validity and cross-study comparability \cite{Deterding2011, Buckley2018};
\item \textbf{M9: Theoretical framework explicitly stated} -- situating findings within explanatory frameworks;
\item \textbf{M10: Study limitations explicitly discussed} -- enabling readers to calibrate confidence in findings.
\end{enumerate}

The MRQS is a measure of \emph{reporting} quality, not of overall study merit \cite{Korniienko2025mrqs}. A qualitative study could score low because several items (M2, M3, M5, M6, M7) are not applicable to its design. This limitation is acknowledged in the interpretation of results.

In addition to the MRQS, risk-of-bias assessment was conducted using an adapted version of the ROBIS tool \cite{Whiting2016}, modified for individual study appraisal. The adapted tool comprised three domains, each with two signalling questions:

\begin{enumerate}
\item \textbf{Domain 1 -- Participant selection and study context}: (Q1) Are the criteria for selecting participants clearly described? (Q2) Is the sample representative of the target population?
\item \textbf{Domain 2 -- Data collection and analysis}: (Q3) Were valid and reliable tools used for data collection? (Q4) Are the methods of data analysis described and justified?
\item \textbf{Domain 3 -- Interpretation and reporting}: (Q5) Are interpretations and conclusions consistent with results? (Q6) Are limitations acknowledged and discussed?
\end{enumerate}

Domain-level judgments followed three rules: (1) both signalling questions ``yes'' $\rightarrow$ low risk; (2) at least one ``no'' $\rightarrow$ high risk; (3) at least one ``unclear'' with no ``no'' $\rightarrow$ unclear risk.

\subsection{Quality assessment results}

\subsubsection{Overall MRQS distribution}

Across the 74 included studies, the mean MRQS was 4.31 out of 10 (SD = 1.93, Mdn = 4.0, IQR = 3--6, range = 0--9) \cite{Korniienko2025mrqs}. No study achieved the maximum score of 10. The distribution was approximately unimodal: 15 studies (20.3\%) scored in the low range (0--2), 38 (51.4\%) in the middle range (3--5), and 21 (28.4\%) in the upper range (6--9). The internal consistency was moderate (Cronbach's $\alpha$ = 0.58), consistent with the instrument's design as a multi-dimensional reporting checklist rather than a single-construct scale.

\subsubsection{Item-level prevalence}

Table~\ref{tab:mrqs-prevalence} presents the prevalence of each MRQS item, revealing three tiers of reporting quality.

\begin{table}[htbp]
\centering
\caption{MRQS item prevalence across 74 studies, with temporal comparison}
\label{tab:mrqs-prevalence}
\small
\begin{tabular}{clcccc}
\hline
\textbf{Item} & \textbf{Criterion} & \textbf{All (n=74)} & \textbf{2004--17} & \textbf{2018--24} & \textbf{$\Delta$} \\
\hline
M9 & Theoretical framework & 70 (94.6\%) & 91.9\% & 97.3\% & +5.4 \\
M1 & Sample size reported & 62 (83.8\%) & 78.4\% & 89.2\% & +10.8 \\
M2 & Sample $\geq$ 30 & 37 (50.0\%) & 51.4\% & 48.6\% & $-$2.7 \\
M10 & Limitations discussed & 35 (47.3\%) & 37.8\% & 56.8\% & +18.9 \\
M6 & Pre-post design & 29 (39.2\%) & 43.2\% & 35.1\% & $-$8.1 \\
M7 & Inferential statistics & 29 (39.2\%) & 37.8\% & 40.5\% & +2.7 \\
M3 & Control group & 23 (31.1\%) & 37.8\% & 24.3\% & $-$13.5 \\
M4 & Validated instruments & 16 (21.6\%) & 18.9\% & 24.3\% & +5.4 \\
M8 & Explicit definition & 11 (14.9\%) & 18.9\% & 10.8\% & $-$8.1 \\
M5 & Effect sizes reported & 7 (9.5\%) & 8.1\% & 10.8\% & +2.7 \\
\hline
\end{tabular}
\end{table}

\textit{High-prevalence items ($>$75\%).} Nearly all studies included a theoretical framework (94.6\%), and most reported sample sizes (83.8\%), indicating that basic framing and descriptive practices are well-established.

\textit{Mid-prevalence items (20--50\%).} Half the studies had sample sizes $\geq$ 30. Limitations discussion showed the greatest improvement over time (+18.9 percentage points). Notably, the use of control groups \emph{declined} from the earlier to the later period ($-$13.5 pp), reflecting the compositional shift toward qualitative and design-based studies.

\textit{Low-prevalence items ($<$20\%).} Only 21.6\% used validated instruments, only 14.9\% provided explicit definitions \cite{30Zin2008, 23Zin2009322, 43ali2018indonesian, 34Serrano2023, 42Huffer201581, 41Petersen2023301, 58Shih2015173, 36Scholz2021, 37Lewis202054, 76Corbeil2011462, 65Razak201139}, and only 9.5\% reported effect sizes. The last finding is particularly consequential: more than 90\% of studies -- including those reporting inferential statistics -- did not report standardised effect measures, severely limiting meta-analytic synthesis \cite{Borenstein2009, Higgins2011, Tamim2011}.

\subsubsection{MRQS by study characteristics}

Journal articles scored slightly higher than conference papers (mean 4.75 vs. 4.17, ns), and studies providing explicit definitions scored higher than those without (mean 5.36 vs. 4.13, p = .065) \cite{Korniienko2025mrqs}. Studies rated as having high overall risk of bias scored lower on the MRQS (mean 4.13) than those rated moderate or lower (mean 5.25).

\subsubsection{Risk-of-bias results}

The overall risk-of-bias assessment revealed \cite{Korniienko2025mrqs, Whiting2016}:
\begin{itemize}
\item High risk of bias: 62 studies (83.8\%)
\item Moderate risk of bias: 9 studies (12.2\%)
\item Low risk of bias: 1 study (1.4\%)
\item Unclear: 2 studies (2.7\%)
\end{itemize}

\textbf{Domain-level analysis:}

\textit{Domain 1 -- Participant selection and study context} was the most persistent source of weakness, with 85.1\% of studies rated as high risk \cite{Korniienko2025mrqs}. The signalling question on sample representativeness (Q2) was the single weakest indicator across the entire instrument: no study received a favourable rating on this criterion. This reflects the universal reliance on convenience sampling (intact classes, available groups) without discussion of how the sample relates to a broader population \cite{Schulz2010CONSORT}.

\textit{Domain 2 -- Data collection and analysis} showed 63.5\% of studies at high risk. The primary concerns were the use of unvalidated instruments (only 21.6\% reported validity or reliability evidence) and insufficient description of analytical procedures.

\textit{Domain 3 -- Interpretation and reporting} was the strongest domain, with 55.4\% at high risk. Studies generally aligned conclusions with results (Q5), but the discussion of limitations (Q6) was met by fewer than half the studies.

\subsection{LLM validation of quality assessment}

The parent systematic review employed LLMs at multiple stages: Claude~3.5 Sonnet for eligibility screening, and GPT-4o for data extraction and risk-of-bias assessment \cite{Korniienko2025mrqs, MalekiKarami2024, Taneri2025, Nyanchoka2020}. To evaluate the reliability of these automated assessments, we conducted a systematic comparison.

\textbf{Risk-of-bias assessment agreement.} Comparison of GPT-4o's automated risk-of-bias responses against human-verified judgments across the six signalling questions for all 74 studies revealed only 37.6\% overall agreement \cite{Korniienko2025mrqs}. Cohen's $\kappa$ values ranged from $-$0.04 to 0.21 across the signalling questions, corresponding to ``poor'' to ``slight'' agreement according to the \citet{LandisKoch1977} benchmarks \cite{Cohen1960, Fleiss1971}. The model consistently applied less stringent criteria than human reviewers, tending to rate studies more favourably on sampling quality (Q1, Q2) and instrument validation (Q3) \cite{Rubinstein2025}.

These low agreement rates contrast markedly with the higher reliability observed for factual data extraction.

\textbf{Factual metadata extraction.} A supplementary multi-LLM comparison involving Gemini~3 Pro, ChatGPT-4, and Grok~4.1-thinking -- each independently extracting data using the same 21-item prompt -- achieved substantially higher agreement \cite{Korniienko2025mrqs}. Pairwise agreement on objective, verifiable fields (sample size, year, publication type, country) ranged from 82\% to 97\%, demonstrating that LLMs are highly reliable for factual extraction but unreliable for subjective quality judgments \cite{MalekiKarami2024, Taneri2025, Rubinstein2025}.

\textbf{Implications for automated research synthesis.} The divergent reliability patterns have important methodological implications \cite{Korniienko2025mrqs}. LLMs can effectively accelerate the labour-intensive process of data extraction in systematic reviews (where they achieve near-human accuracy on factual items) \cite{MalekiKarami2024}. However, they cannot reliably replace human judgment for risk-of-bias assessment, where the evaluative component requires nuanced interpretation of study design features, statistical appropriateness, and reporting completeness \cite{Taneri2025, Rubinstein2025}. The finding that GPT-4o applied less stringent criteria than human reviewers is particularly concerning, as it means that LLM-assisted reviews may systematically underestimate the risk of bias in the included studies.

These findings suggest a hybrid approach: LLMs for initial screening and factual extraction, with human verification for all evaluative judgments. This approach was, in fact, the method employed in the parent systematic review, and its success validates the hybrid model for educational research synthesis.

\subsection{Definitional practices and quality}

The relationship between definitional practices and methodological rigour reveals an important finding \cite{Korniienko2025mrqs}. Studies providing explicit gamification definitions \cite{30Zin2008, 23Zin2009322, 43ali2018indonesian, 34Serrano2023, 42Huffer201581, 41Petersen2023301, 58Shih2015173, 36Scholz2021, 37Lewis202054, 76Corbeil2011462, 65Razak201139} scored higher on the MRQS (mean 5.36 vs. 4.13, p = .065) but showed equivalent risk-of-bias ratings. This suggests that definitional clarity functions as a marker of reporting conscientiousness -- authors who take the time to define their key construct also tend to report other methodological details more thoroughly -- rather than as a direct indicator of design rigour. A study can define gamification clearly while still employing a small sample and no control group.

The decline in definitional clarity from the earlier to the later period ($-$8.1 percentage points) is counterintuitive: one might expect that definitional conventions would strengthen as a field matures. The opposite pattern may reflect the normalisation of gamification terminology -- as the concept becomes familiar, authors may assume that definitions are unnecessary, despite the continued conceptual ambiguity documented in this review.

The inter-item analysis revealed a quantitative-methods cluster (M3 control group, M6 pre-post design, M7 inferential statistics, with phi correlations of 0.36--0.60) operating independently of the conceptual clarity items (M8 definition, M9 framework) \cite{Korniienko2025mrqs}. This structural finding suggests that the methodological quality of gamification research operates along two independent dimensions: \emph{empirical design quality} (related to study design, measurement, and analysis) and \emph{conceptual reporting quality} (related to definitions, frameworks, and limitations). Improving the field requires attention to both dimensions simultaneously \cite{DichevDicheva2017, Clark2016}.

\subsection{Detailed bias analysis by category}

Beyond the aggregate statistics, the risk-of-bias assessment revealed specific patterns of methodological weakness that merit detailed examination.

\subsubsection{Sample size limitations}

A substantial proportion of included studies exhibited significant risk of bias due to inadequate sample sizes \cite{Korniienko2025mrqs}. Of the 62 studies reporting participant numbers, 41 (66.1\%) utilised samples below 50 participants, raising concerns about statistical power and reliability \cite{Cohen1960}. Studies with particularly small samples (n $\leq$ 17) included VR-based studies (n = 17) \cite{1Barbatsis201159}, first-year student evaluations (n = 5) \cite{22Nugroho2020149}, respondent surveys (n = 5), preliminary evaluations (n = 4) \cite{51Ardito2007, 73Wolfartsberger2011160}, and teacher observations with classroom observations (n = 8 teachers). Several studies explicitly acknowledged sample size as a limitation, noting ``limited sample size'' as a constraint or identifying insufficient sample size for robust quantitative analysis.

\subsubsection{Research design concerns}

Significant heterogeneity in research designs contributed to varying levels of bias risk \cite{Korniienko2025mrqs, HigginsThomas2019}. Approximately 25 studies (31\%) lacked comparison groups entirely, relying solely on pre-post single-group designs without controls. This design limitation precludes causal inference regarding the effects of interventions \cite{Schulz2010CONSORT}. Several studies employed purely qualitative methodologies without quantitative outcome measures. While qualitative approaches provide valuable insights into user experience and implementation processes, the absence of standardised outcome metrics introduces subjectivity and limits cross-study comparability \cite{Connolly2012, Clark2016}.

Studies utilising quasi-experimental designs with non-random assignment employed convenience sampling or intact classroom groups, introducing potential selection bias and confounding variables. One study explicitly acknowledged possible instructor bias, as the same teacher delivered both experimental and control conditions, potentially influencing differential treatment effects.

\subsubsection{Measurement and assessment issues}

Numerous studies demonstrated risk of bias related to outcome measurement \cite{Korniienko2025mrqs}. Some relied exclusively on perceived learning rather than objective assessments, utilising Likert scale questionnaires without validation of actual knowledge gains \cite{48Moketar2021, 21Ghulamani2017}. Self-report bias was prevalent, with many studies relying primarily on self-reported data without objective verification \cite{2Mahamarowi2023133, 69Wong2017150}. Instrument validation was frequently absent or inadequately reported: studies used unvalidated survey questionnaires with no reported reliability information, and game testing employed no standardised assessment tools \cite{5kazanidis2018alpha, 20Erandaru2014}.

\subsubsection{Confounding variables and external validity}

Several studies failed to adequately control for potential confounding variables, including prior gaming experience (which could substantially influence engagement and learning outcomes) \cite{49Huizenga2007127, 45Yu2016727} and epistemological beliefs and prior game preferences. The novelty effect represented a recurring threat to validity, where student engagement might be artificially inflated due to the unfamiliarity of game-based approaches rather than their intrinsic educational value \cite{DichevDicheva2017, Hamari2014}.

Generalisability was severely limited in numerous studies. Single-institution studies could not account for variations in educational systems, cultural contexts, or institutional resources that might moderate intervention effects. Geographic and cultural constraints further limited external validity: the concentration of studies in a small number of countries (Taiwan, USA, Malaysia, Indonesia accounting for 55.4\%) means that findings may not generalise to other educational contexts.

\subsubsection{Implementation fidelity}

Implementation fidelity issues introduced additional bias risks \cite{Korniienko2025mrqs}. Studies reported significant technical difficulties including GPS/connectivity problems, hardware limitations, and AR recognition issues that interfered with intervention delivery \cite{58Shih2015173, 51Ardito2007, 25Repantis2011502}. Teacher-related factors represented another concern, with variability in implementation approaches: some teachers used games merely as visual aids rather than engaging students in active historical inquiry, potentially diluting intervention effects \cite{61Karsenti202027}. Time constraints affected multiple studies, with insufficient class time for adequate implementation leading to incomplete student experiences \cite{49Huizenga2007127, 46Stirling2022}.

\subsection{Implications for the field}

The high prevalence of methodological weaknesses has several important implications:

First, the positive findings regarding gamification effectiveness reported in Section~\ref{sec:gamification-review} should be interpreted cautiously. Studies with high risk of bias may overestimate effect sizes. Previous meta-analyses in educational technology have shown that studies with higher methodological quality tend to report smaller effect sizes \citep{Hattie2009}, suggesting that the true effectiveness of gamification in history education may be more modest than the current literature implies.

Second, the field urgently needs higher-quality research designs \cite{Korniienko2025mrqs, Schulz2010CONSORT}. The decline in control group usage ($-$13.5 pp from 2004--2017 to 2018--2024) is particularly concerning, as it means the newer literature is less capable of supporting causal inference than the older literature. Quasi-experimental designs with adequate sample sizes ($\geq$ 64 per group for medium effects) and pre-registered protocols should become the standard for gamification effectiveness research \cite{GuyattGRADE2008}.

Third, the definitional clarity issue is not merely a semantic concern but a methodological one \cite{Deterding2011, Buckley2018}. Without clear definitions, construct validity cannot be established, and it becomes impossible to determine whether different studies are investigating the same phenomenon \cite{DichevDicheva2017}. The field would benefit from adopting a shared taxonomy that distinguishes structural gamification, content gamification, full gamification, and game-based learning \cite{36Scholz2021, SEABORN201514}.

Fourth, the near-universal failure to report effect sizes (90.5\% of studies) prevents meta-analytic synthesis and makes it impossible to estimate the practical significance of gamification effects \cite{Borenstein2009, Higgins2011}. Future studies should routinely report Cohen's d, eta-squared, or equivalent measures alongside inferential test results \cite{Cohen1960}.

Fifth, the LLM validation findings have important implications for the growing practice of AI-assisted systematic reviews \cite{MalekiKarami2024, Taneri2025}. While LLMs can accelerate factual data extraction (82--97\% agreement), they cannot replace human judgment for evaluative tasks like risk-of-bias assessment ($\kappa$ = $-$0.04 to 0.21) \cite{LandisKoch1977}. Reviews that rely on LLM-only quality assessment may systematically underestimate methodological problems \cite{Rubinstein2025}.

\subsection{Proposed reporting checklist}

Based on the MRQS analysis \cite{Korniienko2025mrqs}, we propose that future studies on gamification in history education should satisfy the following minimum reporting requirements:

\begin{enumerate}
\item Provide an explicit definition of gamification (or game-based learning, serious games, etc.) and position the study within the four-category taxonomy (structural, content, full, game-based learning) \cite{Deterding2011, 36Scholz2021, Buckley2018};
\item Report the sample size, sampling method, and any available demographic characteristics \cite{Schulz2010CONSORT};
\item Describe the measurement instruments used and provide evidence of validity and/or reliability;
\item Report effect sizes alongside inferential statistics \cite{Borenstein2009, Cohen1960};
\item Include a comparison condition (even if non-equivalent) where the research question involves effectiveness claims \cite{Schulz2010CONSORT, HigginsThomas2019};
\item Discuss study limitations explicitly, including threats to internal and external validity;
\item Specify the theoretical framework guiding the design and interpretation;
\item Report the duration of the intervention and the timing of assessments;
\item Describe the game design elements employed and the rationale for their selection \cite{SEABORN201514, Sailer2017};
\item Make data and materials available where possible to support replication.
\end{enumerate}

This checklist operationalises the MRQS items into actionable guidance for researchers. If adopted, it would address the most critical reporting gaps identified in the current literature and move the field toward the transparent, rigorous evidence base that practitioners need.

\subsection{Conclusions to Section 1.3}

The MRQS quality assessment \cite{Korniienko2025mrqs} reveals that the vast majority (83.8\%) of gamification studies in history education carry a high risk of bias. The mean MRQS was 4.31/10, with the most critical weaknesses being the near-universal failure to report effect sizes (90.5\%) \cite{Borenstein2009}, the rarity of validated instruments (21.6\%), and the pervasive definitional ambiguity (85.1\%) \cite{Deterding2011, Buckley2018}. Domain-level risk-of-bias analysis \cite{Whiting2016} identified participant selection (Domain~1, 85.1\% high risk) as the most persistent source of weakness, with no study meeting the representativeness criterion. The LLM validation study revealed that automated risk-of-bias assessment achieves only slight agreement with human reviewers ($\kappa$ = $-$0.04 to 0.21) \cite{LandisKoch1977, Fleiss1971}, while factual data extraction is highly reliable (82--97\% agreement) \cite{MalekiKarami2024, Taneri2025, Rubinstein2025}. These findings collectively inform both the design of the experimental study (Chapter~\ref{ch:methodology}) and the self-assessment approach adopted in the discussion (Chapter~\ref{ch:results}).


\section{Gamification and digital competence: implementation through DigComp 2.0} \label{sec:digcomp}

\subsection{The DigComp 2.0 framework}

The European Digital Competence Framework for Citizens (DigComp 2.0) provides a comprehensive reference framework for digital competence \citep{Vuorikari2016}. The framework identifies five competence areas:

\begin{enumerate}
\item \textbf{Information and data literacy}: Browsing, searching, and filtering data; evaluating data; managing data.
\item \textbf{Communication and collaboration}: Interacting through digital technologies; sharing through digital technologies; engaging in citizenship through digital technologies; collaborating through digital technologies; netiquette; managing digital identity.
\item \textbf{Digital content creation}: Developing digital content; integrating and re-elaborating digital content; copyright and licences; programming.
\item \textbf{Safety}: Protecting devices; protecting personal data and privacy; protecting health and well-being; protecting the environment.
\item \textbf{Problem solving}: Solving technical problems; identifying needs and technological responses; creatively using digital technologies; identifying digital competence gaps.
\end{enumerate}

\subsection{Mapping gamification to DigComp 2.0}

The implementation of gamification in history education can be systematically aligned with DigComp 2.0 competence areas \citep{Titova2025cte} \cite{Korniienko2025cte}. Table~\ref{tab:digcomp-mapping} presents a detailed mapping of gamification activities to specific DigComp sub-competencies, demonstrating that gamified history education simultaneously develops historical knowledge and digital competencies \cite{Voogt20161}.

\begin{table}[htbp]
\centering
\caption{Mapping of gamification activities to DigComp 2.0 competence areas}
\label{tab:digcomp-mapping}
\small
\begin{tabular}{p{3cm}p{4cm}p{5.5cm}}
\hline
\textbf{DigComp area} & \textbf{Sub-competencies} & \textbf{Gamification activities} \\
\hline
1. Information and data literacy & 1.1 Browsing, searching; 1.2 Evaluating; 1.3 Managing & Navigating Google Forms quizzes; evaluating historical sources in quiz questions; managing responses across 156 forms \\
\hline
2. Communication and collaboration & 2.1 Interacting; 2.2 Sharing; 2.4 Collaborating; 2.5 Netiquette & Viber group interactions; sharing resources and quiz results; peer support; appropriate communication norms \\
\hline
3. Digital content creation & 3.1 Developing; 3.2 Integrating & Interacting with HistoryLikeGame platform; engaging with multimedia course materials; creating responses in varied digital formats \\
\hline
4. Safety & 4.1 Protecting devices; 4.3 Health and well-being & Wartime digital safety awareness; balancing online learning with physical safety; managing screen time \\
\hline
5. Problem solving & 5.1 Technical problems; 5.2 Identifying needs; 5.3 Creative use & Solving quiz problems; adapting to technical disruptions; finding alternative access methods during outages \\
\hline
\end{tabular}
\end{table}

\textbf{Information and data literacy (Area~1).} History quiz-based gamification directly engages students with information processing tasks: interpreting primary sources presented in digital format, evaluating the reliability of historical evidence, and managing data across multiple assessment forms \cite{Vuorikari2016, Korniienko2025cte}. In the context of this study's 156 Google Forms quizzes, students practised data literacy with each assessment interaction. The multi-block structure of seminars (typically 3--7 quiz blocks per session) required students to navigate between forms, track their progress, and manage their responses -- activities that develop competencies in data management (sub-competency 1.3).

\textbf{Communication and collaboration (Area~2).} The Viber chat group used in the course served dual purposes: facilitating communication about course logistics and creating a space for collaborative learning \cite{Korniienko2025cte, Vuorikari2016}. Students shared resources, discussed historical questions, and provided peer support. The gamified naming conventions for quiz blocks (using themes from Greek mythology, gemstones, flowers, and minerals) created a shared vocabulary that fostered community and demonstrated the social function of creative digital communication \cite{symbolismhub.com, Guitert2023245}. The chat log analysis documented 400 messages reflecting various communication patterns: formal absence notifications, informal peer interaction, resource sharing, and instructor support -- all contributing to sub-competency 2.5 (netiquette) and 2.4 (collaborating through digital technologies).

\textbf{Digital content creation (Area~3).} The gamified platform HistoryLikeGame engaged students with interactive digital content including drag-and-drop matching exercises and text-input challenges \cite{Korniienko2025cte}. While students were primarily consumers rather than creators of digital content in this study, the platform's HTML-based architecture demonstrated how educational digital content can be created and deployed using standard web technologies \cite{Vuorikari2016}. The open-source nature of the platform (hosted on GitHub Pages) provides a model for educators interested in developing their own gamification tools, addressing sub-competency 3.1 (developing digital content).

\textbf{Safety and well-being (Area~4).} The wartime context of this study added a critical and unique dimension to digital safety that is absent from most DigComp implementations \cite{Korniienko2025cte, Vuorikari2016}. Students navigated online learning while managing real threats including power outages, internet disruptions, and air raid alerts. The course's asynchronous components and archived materials ensured continued access to education even when synchronous participation was impossible. The instructor's explicit safety guidance (``Your safety is the most important [...] if in doubt, go to a shelter'') modelled responsible digital behaviour that prioritises physical well-being over digital participation -- a competency area (4.3: protecting health and well-being) that takes on heightened significance in conflict zones \cite{UNESCO2023}.

\textbf{Problem solving (Area~5).} Gamified assessments inherently involve problem-solving: students must identify relevant information, apply historical knowledge to novel contexts, and select appropriate strategies for different question types (matching, sequencing, multiple choice, open-ended) \cite{Vuorikari2016, Korniienko2025cte, Osuna-Acedo202189}. The wartime context further developed problem-solving competencies (sub-competency 5.1: solving technical problems) as students adapted to infrastructure disruptions by finding alternative ways to access course materials and complete assessments.

\subsection{Implementation guidelines}

Based on the synthesis of the systematic reviews and the DigComp 2.0 alignment analysis \cite{Korniienko2025cte, Vuorikari2016}, the following guidelines for implementing gamification in history education were developed. Each guideline is grounded in the evidence base established in Sections~\ref{sec:dh-review}--\ref{sec:mrqs} and is connected to specific DigComp competence areas.

\begin{enumerate}
\item \textbf{Start with clear learning objectives} (DigComp Areas 1, 5): Define specific historical knowledge outcomes and digital competencies to be developed before selecting gamification elements. The MRQS analysis showed that studies with explicit theoretical frameworks (94.6\%) generally produced more interpretable results. The learning objectives should specify: (a) the historical content to be mastered, (b) the historical thinking skills to be developed, and (c) the digital competencies to be practised.

\item \textbf{Use progressive complexity} (DigComp Area 5): Structure gamified activities in a sequence of increasing difficulty, from recognition tasks (matching, multiple choice) through analysis tasks (source evaluation, chronological sequencing) to synthesis tasks (constructing narratives, arguing causation). This progression mirrors the cognitive development sequence identified in the design principles (Section~\ref{sec:gamification-review}) and aligns with established taxonomies of historical thinking.

\item \textbf{Maintain historical rigour} (DigComp Area 1): Ensure that game elements do not trivialise or distort historical content. The systematic review found that 10.8\% of studies reported negative outcomes, often associated with superficial gamification that prioritised engagement over accuracy. All quiz questions should be based on verified historical facts and aligned with curriculum standards (in this study, the NMT Ukrainian History syllabus).

\item \textbf{Provide multiple pathways} (DigComp Areas 3, 5): Offer both required and optional (``extra'') gamified activities to accommodate different ability levels and engagement preferences. The HistoryLikeGame platform exemplifies this approach with its five sections (plus Omega comprehensive review), allowing students to choose their level of engagement. Wordwall and other supplementary platforms provide additional variety.

\item \textbf{Integrate assessment with learning} (DigComp Areas 1, 5): Use gamified quizzes as both formative assessment tools and learning activities, providing immediate feedback and opportunities for retry. The case of Talash D. (Chapter~\ref{ch:results}), who submitted 7 attempts at a comprehensive test with progressively improving scores, demonstrates the power of this approach for iterative mastery learning.

\item \textbf{Ensure accessibility and resilience} (DigComp Area 4): Design for various devices and connectivity conditions, particularly important in contexts where infrastructure may be unreliable. The HistoryLikeGame platform's static HTML architecture (total size: 205~KB) provides a model for accessible design. In wartime or disaster-affected contexts, the five-layer resilience model (synchronous, communication, content, assessment, gamification) ensures educational continuity.

\item \textbf{Build community through gamification} (DigComp Area 2): Use creative naming, shared references, and collaborative elements to foster a sense of belonging and shared purpose among learners. The thematic quiz naming conventions (gemstones, Greek mythology, flowers, trees, minerals) served this function in the present study, creating cultural enrichment opportunities alongside the history content.

\item \textbf{Document systematically} (DigComp Areas 1, 3): Maintain detailed records of implementation (chat logs, assessment data, platform analytics) to enable evidence-based evaluation and future research. The MRQS analysis revealed that comprehensive documentation is the exception rather than the rule in gamification research -- only 47.3\% of studies discussed limitations, and only 9.5\% reported effect sizes. Systematic documentation addresses the most critical reporting gaps identified in Section~\ref{sec:mrqs}.
\end{enumerate}

These guidelines synthesise theoretical insights from the four analyses conducted in this chapter (digital history review, gamification review, MRQS assessment, and DigComp alignment) into actionable recommendations for practitioners. Their application in the experimental study is documented in Chapter~\ref{ch:methodology}.

\subsection{Conclusions to Section 1.4}

The alignment of gamification with DigComp 2.0 \cite{Vuorikari2016} demonstrates that gamified history education can simultaneously develop historical knowledge and all five areas of digital competence: information and data literacy, communication and collaboration, digital content creation, safety, and problem solving \cite{Korniienko2025cte, Titova2025cte}. The wartime context of this study adds a unique dimension to the safety competence area, as students must navigate real physical threats alongside digital learning challenges \cite{UNESCO2023}.

The eight implementation guidelines derived from the synthesis of systematic reviews \cite{KorniienkoSemerikov2024, Korniienko2025mrqs, Korniienko2025els} provide a practical framework for educators seeking to implement gamification in their history courses. These guidelines address the key weaknesses identified in the MRQS assessment (lack of clear objectives, insufficient documentation, absent definitions) \cite{Korniienko2025mrqs} while operationalising the six design principles from the gamification review (historical authenticity, narrative integration, progressive difficulty, immediate feedback, social elements, multiple modalities) \cite{KorniienkoSemerikov2024}.


\section{Conclusions to Chapter 1} \label{sec:ch1-conclusions}

The theoretical analysis presented in this chapter establishes the following key findings:

\begin{enumerate}
\item \textbf{Digital history pedagogy is an emerging field}: The PRISMA-guided systematic review \cite{Page2021, Korniienko2025els} identified only 6 pedagogical studies \cite{Bell2016, Coleborne2011, Davis2017, Lee2014159, McLean20144, Soh2013}, indicating that digital history in education lacks the critical mass of research needed for evidence-based practice (Section~\ref{sec:dh-review}).

\item \textbf{Gamification is the most researched digital approach to history education}: The systematic review of 74 studies \cite{KorniienkoSemerikov2024} confirms that gamification has attracted the most research attention among digital history methods, but the field is characterised by definitional ambiguity (85.1\% lack definitions) \cite{Deterding2011} and predominantly positive but methodologically weak evidence (Section~\ref{sec:gamification-review}).

\item \textbf{Methodological quality is critically low}: The MRQS assessment \cite{Korniienko2025mrqs} reveals that 83.8\% of gamification studies in history education carry a high risk of bias, primarily due to small samples, lack of randomisation, and insufficient reporting. The LLM validation approach \cite{MalekiKarami2024, Taneri2025, Rubinstein2025} offers a promising pathway for scaling quality assessment (Section~\ref{sec:mrqs}).

\item \textbf{Gamification can develop digital competencies}: Systematic alignment with DigComp 2.0 \cite{Vuorikari2016} demonstrates that gamified history education activities address all five competence areas \cite{Korniienko2025cte, Titova2025cte}, making gamification not merely an engagement tool but a vehicle for comprehensive digital competence development (Section~\ref{sec:digcomp}).

\item \textbf{Evidence-based guidelines are needed}: The synthesis of findings from all four analyses \cite{Korniienko2025els, KorniienkoSemerikov2024, Korniienko2025mrqs, Korniienko2025cte} provides the foundation for implementation guidelines that address the identified gaps: clear definitions \cite{Deterding2011, Buckley2018}, rigorous design \cite{Schulz2010CONSORT}, progressive complexity, and systematic documentation.
\end{enumerate}

These findings collectively justify the experimental study presented in the following chapters: a gamification-enhanced digital history course designed according to the identified principles \cite{KorniienkoSemerikov2024, Korniienko2025mrqs}, implemented in a real educational setting, and documented with sufficient detail to address the methodological weaknesses prevalent in the field \cite{Korniienko2025mrqs, Schulz2010CONSORT}.
